\section{Model, noise laws, and output contracts}\label{sec:model}

This section fixes the input model, the perturbation laws, the meaning of
an exact smoothed algorithm, and the output formats. Later sections state
their theorems in these terms. The section ends with two elementary facts
that apply to all results: a certified bound for the original objective, and
an example showing why a finite law forces some exact treatment of ties.

\subsection{Instances and input length}\label{sec:model:input}

All data are rational and encoded in binary. A \emph{base instance}
consists of an objective $F_0$, a feasible set $X\subseteq\R^n$, a noise
scale $\sigma>0$, and the structural data required by the relevant theorem:
a quadratic convexifier, a tree decomposition, a core, curvature bounds, or
certificates. Its \emph{base length} $I$ is the total binary length of these
data. It is measured before any noise is drawn; the bit lengths of sampled
coefficients are accounted for separately.

\paragraph{Objectives.}
A quadratic objective is $F_0(x)=\frac12x^{\mathsf T}Ax+b^{\mathsf T}x+c$
with a symmetric rational matrix $A$. A polynomial objective of total degree
at most $d$ is given by an explicit list of monomials, each a rational
coefficient together with the list of its variables and their positive
exponents, or as a sum of factors, each such a list on a stated set of
variables (its \emph{scope}). This sparse encoding charges a monomial for the
variables it contains, not for the whole dimension. In every polynomial theorem
$d$ is a fixed constant, and $\poly_d(\cdot)$ denotes a polynomial whose
exponent depends on $d$ only. Objectives are never given as arithmetic
circuits. Circuits can be exponentially more compact than monomial lists,
and they change the complexity of evaluation and of exact output;
\cref{sec:lim:points} uses this difference.

\paragraph{Domains.}
A \emph{mixed box} is $X=\prod_{i=1}^nX_i$, where each $X_i$ is either a
continuous interval $[\ell_i,u_i]$ or a \emph{native integer} interval
$[\ell_i,u_i]\cap\Z$. Native integer bounds are encoded in binary, so the
number of labels $u_i-\ell_i+1$ can be exponential in $I$. Integer bounds
are rounded inward, empty domains are rejected, and fixed coordinates are
substituted before any algorithm starts. After this preprocessing every
remaining coordinate has width $w_i=u_i-\ell_i>0$, and the
\emph{continuous hull} of $X$ is $\prod_i[\ell_i,u_i]$. We write $n_c$ and
$n_z$ for the numbers of continuous and native integer coordinates. A
\emph{mixed polytope} is $P\cap(\Z^{n_z}\times\R^{n_c})$ for a bounded
rational polyhedron $P$ given by linear inequalities. Products of simplices,
order polytopes, explicit and implicit graph constraints, integer flows, and
totally unimodular systems are defined where they are used.
Optimization statements assume $X\ne\varnothing$, or explicitly decide
feasibility and report an empty set before promising an attaining optimizer.

\paragraph{Structural and numerical parameters.}
The structural parameters are the negative inertia $k=n_-(A)$ of a quadratic
objective, the number $k$ of rows in a supplied concave factor when that
factor is used, the bag size $p$ (treewidth plus one) of a supplied tree
decomposition whose bags contain all factor scopes, the size $k$ of a
supplied continuous core, and, in the integer theorems, the integer
dimension. The numerical parameters are curvature and width quantities
compared with $\sigma$. Examples are an upper bound $L$ on coordinate second
derivatives, $\nu=\max\{0,-\lambda_{\min}(A)\}$, a convexifier scale
$\alpha$, and coordinate widths $w_i$. Each theorem states on which set a
curvature bound must hold. A numerical parameter enters a bound through its
magnitude, not through its encoding length.

\paragraph{Status of premises.}
Each premise of a theorem has one of three statuses. It is \emph{verified}
if the algorithm decides it within the stated work bound; examples are exact
definiteness tests of rational matrices and a curvature bound computed from
monomial coefficients. It is \emph{certified} if the input contains a proof
of it: the proof's encoding length is part of $I$, and the work of checking
it is part of the stated work bound. It is \emph{promised} if the algorithm
uses it without checking. A theorem says, for each promise, whether
correctness depends on it or only the complexity analysis does. A false
promise of the first kind can produce a wrong output; a false promise of the
second kind leaves every output correct and can only invalidate the stated
work bound.

\subsection{Linear perturbations and finite laws}\label{sec:model:noise}

Every result perturbs only the linear part of the objective. The
\emph{sampled objective} is
\begin{equation}\label{eq:model:sampled}
 F_\gamma(x)=F_0(x)+\gamma^{\mathsf T}x,
\end{equation}
for a random vector $\gamma\in\Q^n$ drawn once. The feasible set, the
nonlinear coefficients, and all structural data are not perturbed.

\begin{definition}[Grid law]\label{def:model:grid}
For a rational $\sigma>0$ and a power of two $M\ge2$, let
\[
 \Lambda_{\sigma,M}=\Bigl\{-\sigma+\frac{2\sigma j}{M-1}:j=0,1,\ldots,M-1\Bigr\}
\]
be the $M$ equally spaced points of $[-\sigma,\sigma]$, including both
endpoints, and let $U_{\sigma,M}$ be the uniform distribution on
$\Lambda_{\sigma,M}$.
\end{definition}

A draw from $U_{\sigma,M}$ uses exactly $\log_2M$ fair random bits, and every
atom has binary length $O(\log M+\operatorname{size}(\sigma))$. Because $M-1$
is odd, $0$ is not an atom. The atoms are equally spaced at distance
$2\sigma/(M-1)$, so a closed interval of length $\lambda$ contains at most
$\lambda(M-1)/(2\sigma)+1$ of them. Hence
\begin{equation}\label{eq:model:interval}
 U_{\sigma,M}(J)\le\frac{\lambda}{2\sigma}+\frac1M
 \qquad\text{for every closed interval $J$ of length $\lambda$.}
\end{equation}
The first term is the bound for the continuous uniform law on
$[-\sigma,\sigma]$. The second is an \emph{atomic term}; it does not
decrease with $\lambda$, and every finite-law estimate in this paper carries
a term of this kind.

\begin{definition}[Perturbation models]\label{def:model:perturbation}
Let $\mathcal D$ be a finite rational law on $\Q$. The following models
specify $\gamma$ in \eqref{eq:model:sampled}.
\begin{enumerate}[label=(\roman*)]
\item \emph{Ambient}: $\gamma_1,\ldots,\gamma_n$ are independent with law
$\mathcal D$.
\item \emph{Core-only}: for a supplied core $C\subseteq[n]$, the coefficients
$\gamma_i$, $i\in C$, are independent with law $\mathcal D$, and
$\gamma_i=0$ for $i\notin C$.
\item \emph{Aligned}: for a supplied rational matrix $T\in\Q^{k\times n}$,
$\gamma=T^{\mathsf T}\xi$, where $\xi_1,\ldots,\xi_k$ are independent with
law $\mathcal D$.
\end{enumerate}
\end{definition}

The models coincide in special cases: core-only noise with $C=[n]$, and
aligned noise with $T=I$, are ambient noise. In general they differ, and a
guarantee proved for one model does not transfer automatically to another.
Aligned coefficients are dependent in the original coordinates and lie in
the row space of $T$. Core-only noise leaves every residual cost unchanged,
so ties and flat directions among residual variables can persist on every
draw.

Two marginal laws are used. The first is the grid law $U_{\sigma,M}$. The
second is a \emph{Gaussian-like} finite rational law $\mathcal G_{b,\sigma}$,
whose Kolmogorov distance to the normal law $N(0,\sigma^2)$ is at most
$2^{-b}$ and which is produced by a bounded rejection sampler from
polynomially many fair bits (\cref{sec:count:laws}). We call $\sigma$ the
\emph{noise scale}: it is the half-width of the grid law and the standard
deviation of the normal proxy. The support of $\mathcal G_{b,\sigma}$ is
contained in $[-(b+20)\sigma,(b+20)\sigma]$. The
exact normal law is used only as a proof device; it is not an input model,
because its samples are not finite rational data.
Our samplers have a bound on their worst-case number of random bits, which
forces finite support. A Turing machine without such a bound can also sample
a rational law with countably infinite support.

\paragraph{Base-chosen laws.}
In every theorem except the strong-noise theorem, the law is computed
deterministically from the base instance before sampling: the integer $M$,
or the accuracy $b$ of the Gaussian-like law, is a function of the base data.
As stated, the law therefore depends on the instance. The strong-noise theorem
(\cref{thm:int:strong-field}) holds for every grid size $M$, and its work
bound depends on $M$ explicitly. The theorems assert their bounds for this
constructed law. They do not assert the same bounds for every finite law
with noise scale $\sigma$, for a coarse grid chosen independently of the
instance, or for an arbitrary approximation of a Gaussian. The reason is the
atomic term in \eqref{eq:model:interval}: the law must be fine enough that
atoms on exceptional sets have small total probability, and ``fine enough''
is measured against base-only quantities, such as the cost of an exact
fallback or the spacing of a value lattice. In the results with polynomial
sampling precision the logarithms of these quantities are bounded by
polynomials in the input length alone, so the law can also be chosen from
the input length rather than from the individual instance
(\cref{prop:model:uniform}).

\paragraph{Sampling precision.}
Write $b=\log_2M$ for the number of random bits per perturbed coordinate of a
grid law. A theorem has \emph{polynomial sampling precision} if
$b\le\poly(I)$, so that the sampled instance has length
$I+n\cdot\poly(I)$. A few results for nonlinear flows and totally
unimodular systems need \emph{parameter-dependent sampling precision}
$b\le f(k)\poly(I)$; this enlarges the sampled instance, and the stated work
bounds account for it.

\begin{proposition}[Resolution depending only on the input length]\label{prop:model:uniform}
Fix one of the theorems of \cref{sec:qp,sec:sp,sec:con,sec:rec,sec:int}
with polynomial sampling precision, together with its fixed format
constants, such as the degree and the dependency depth of graph constraints,
and its oracle implementation. There are effective nondecreasing
positive integer-valued polynomials $P$, $A$, and $D$, depending only on these choices, such that the
theorem remains valid, with the same output contract and the same
structural and numerical factors in its expected work bound up to a change
of its polynomial factor, when its law is replaced as follows for every
valid base instance of length at most $I$:
\begin{enumerate}[label=(\alph*)]
\item a grid law by $U_{\sigma,M(I)}$ with $M(I)=2^{P(I)}$; in the lattice
theorem (\cref{thm:int:lowrank}) the level cap is reset to $\log_2M(I)$;
\item a Gaussian-like law by $\mathcal G_{b(I),\sigma}$ with
$b(I)=A(I)+D(I)t(I)$ and $t(I)=\lceil4\log_2(A(I)+D(I)+21)\rceil$, the
auxiliary box and the level cap being recomputed for this support.
\end{enumerate}
For the core-only flow and totally unimodular theorems with arbitrary core
faces (\cref{thm:int:flow-boundary,thm:int:tu}) the same holds for instances
whose core size is at most a supplied $K$, with $M(I,K)=2^{\lceil
F(K)P(I)\rceil}$ for an effective nondecreasing $F$, and with $F(K)$ in the
sampling, output, and work bounds. The integer envelope $F(K)$ is chosen
with evaluation cost polynomial in $1+\operatorname{size}(K)+F(K)$,
and this computation is charged before sampling. The strong-noise theorem holds for every
$M$; its probabilities $q_i$ must be computed for the chosen $M$, and its
sufficient regime is preserved by every larger power of two.
\end{proposition}

\begin{proof}
Apply \cref{cor:count:universal-law}, whose routewise verification is in
\cref{app:count:universal}. The effective bounds control logarithms of the
base-only format and degree multipliers; in particular
$\log(s+1)\le\poly_d(I)$ for an atom universe of size $s$, even when native
integer ranges contribute exponentially many atoms. Coefficient heights and
fallback work retain their separate polynomial dependence on the added
sampled bits. The Gaussian support schedule recomputes the auxiliary box and
level cap, and gives $(b(I)+20)\sigma\le2^{t(I)}\sigma$.
The lattice cap is reset as stated. The arbitrary-face flow and TU cases use
the parameter-dependent part of the corollary with $k\le K$.
For strong noise the actual $q_i$ and $\beta$ are recomputed; the sufficient
regime of \cref{eq:int:sf-regime} holds for every larger power of two.
\end{proof}

The resolution $M(I)$ standardizes the scalar law, which is still rescaled
by the supplied $\sigma$ and applied to the coordinates that the theorem
perturbs; it is not one probability measure for all instances. Taking
$K=I$ in the flow case gives a law depending only on $I$, but its sampling
cost need not stay within $f_d(k)\poly_d(I)$ for the actual core size $k$.
The geometric and oracle qualifications remain part of the proposition.
The routewise widths and derivative bounds have polynomial encoding length;
fixed degree alone does not ensure this for arbitrary compact semialgebraic
domains. For example, $1\le x_0\le2$, $x_i=x_{i-1}^2$ has last-coordinate
width $2^{2^N}-1$. The proposition also fixes the oracle implementation; it
does not supply one for a broader class.

\paragraph{Small and strong noise.}
In most results $\sigma$ is a free parameter, and the bounds degrade as
$\sigma$ decreases; these are statements about small perturbations of an
arbitrary instance. The random-component results of \cref{sec:int} instead
require noise that dominates the variation of every coordinate derivative.
There $\sigma$ is large compared with the curvature of the instance, and the
results describe instances whose linear part is strongly random rather than
slightly perturbed.

\subsection{Exact smoothed algorithms}\label{sec:model:algorithms}

Computation is measured in bit operations of a deterministic Turing
machine. Oracle calls, where a theorem allows them, are charged at their
stated cost.

\begin{definition}[Exact smoothed algorithm]\label{def:model:algorithm}
Fix a class of base instances, an output format, a perturbation model, and a
rule assigning to each base instance $\Pi$ a finite rational law
$\mathcal D_\Pi$ for the perturbation. An \emph{exact smoothed algorithm with
expected work $W$} is a deterministic algorithm $\mathcal A$ with the
following properties.
\begin{enumerate}[label=(\roman*)]
\item\label{it:model:every} For every base instance $\Pi$ and every $\gamma$
in the support of $\mathcal D_\Pi$, $\mathcal A(\Pi,\gamma)$ halts and
returns an exact global optimizer and the optimal value of $F_\gamma$ on $X$,
in the stated format.
\item\label{it:model:once} The perturbation is drawn once. The algorithm
does not resample, discard, or condition on any event of the draw.
\item\label{it:model:expect} The bit work of computing $\mathcal D_\Pi$,
sampling, and running $\mathcal A(\Pi,\gamma)$ has expectation at most
$W(\Pi)$ over $\gamma\sim\mathcal D_\Pi$.
\end{enumerate}
\end{definition}

The randomness defines the instance to be solved; the algorithm itself is
deterministic given the draw. Property~\ref{it:model:every} includes draws
with tied or positive-dimensional optimal sets and draws on which a fast
route fails. Property~\ref{it:model:once} excludes the shortcut of redrawing
until a draw is generic, which would change the law of the solved instance
and with it the meaning of the expectation. By Markov's inequality, an
expected bound $W$ implies work at most $tW$ with probability at least
$1-1/t$; it is not a worst-case bound, and on rare draws the work can be
exponential in $I$.

\paragraph{How ties are handled.}
A finite law gives ties and other degenerate draws positive probability
(\cref{sec:model:atom}), so an exact smoothed algorithm needs a way to
solve them. Most results use a fast route that refines a search and closes
it with an exact certificate, together with an exact \emph{fallback}: a
deterministic algorithm that solves every sampled instance with bit work at
most $B\cdot\poly(I+b)$, where the factor $B$ depends only on the base
instance and $\log B\le\poly(I)$. The fallback runs on the same draw when
the fast route has not finished by a stopping depth computed from base data.
If this happens with probability at most $1/B$, the fallback contributes
only a polynomial term in $I+b$ to the expected work, including the stated
parameter factor when sampling precision depends on the core size
(\cref{lem:count:rare-fallback}). The quantities are chosen in a fixed order:
first $B$, then the thresholds that define exceptional events, then the
stopping depth, and finally the law. Two results need no rare fallback. The
lattice theorem for low-rank integer problems (\cref{thm:int:lowrank}) stops
exactly once the certified gap falls below the spacing of the lattice of
attainable objective values. The strong-noise component theorem
(\cref{thm:int:strong-field}) solves every random component exactly on every
draw.

\paragraph{Expectations of adaptive searches.}
The fast routes refine cells adaptively, and which cells are visited
depends on the draw. Expected counts are therefore never computed as if the
retained cells were independent of the noise. Every visited cell is charged
to a node of a deterministic grid fixed before sampling, and the expected
number of nodes that satisfy a necessary near-optimality condition is
bounded node by node (\cref{thm:count:cells}).

\subsection{Output contracts}\label{sec:model:outputs}

Exact output for $F_\gamma$ takes one of the following forms. Each theorem
states which form it returns on ordinary draws and on fallback draws.

\begin{definition}[Output formats]\label{def:model:outputs}
\leavevmode
\begin{enumerate}[label=(\alph*)]
\item\label{it:model:rational} \emph{Rational output}: a feasible
$x^*\in\Q^n$ with $F_\gamma(x^*)=\min_XF_\gamma$, and this value.
\item\label{it:model:implicit} \emph{Implicit patch output}: a set $S$ of
coordinates with fixed rational values $z_S$ (native integer labels,
original continuous bounds, or continuous values fixed by a sound
singleton-hull certificate), a rational box $Q$ in the remaining
coordinates such that $(z_S,y)\in X$ for every $y\in Q$, and a rational
$g_0>0$ with a verified proof that
$\nabla^2_{yy}F_\gamma(z_S,y)\succeq g_0I$ on $Q$, together with the claim
that the unique minimizer $\hat y$ of $F_\gamma(z_S,\cdot)$ on $Q$ gives a
global optimizer $\hat x=(z_S,\hat y)$ of $F_\gamma$ on $X$. Products of
simplices and order polytopes use a rational relative polytope in place of
$Q$.
\item\label{it:model:charted} \emph{Charted output}: for feasible sets that
are graphs of maps from free to dependent coordinates, an implicit patch
output for the reduced objective in the free coordinates, together with a
certified chart that determines every dependent coordinate from the free
ones. The chart is either an explicit polynomial map, or a list of scalar
equations, each with a rational bracket on which a certified positive
derivative bound guarantees a unique root. The optimizer is the lift of the
unique reduced patch minimizer through the chart.
\item\label{it:model:algebraic} \emph{Algebraic output}: a real algebraic
number $\theta$, given by a nonzero integer polynomial and an isolating
interval with rational endpoints, and rational polynomial maps that give
every continuous coordinate of one optimizer and the optimal value as
functions of $\theta$, with integer coordinates listed explicitly.
\item\label{it:model:component} \emph{Component output}: a partition of the
unfixed coordinates into components, an algebraic output for each
component, rational values for the remaining coordinates, and the optimal
value written as a rational constant plus the sum of the component values.
\end{enumerate}
\end{definition}

Several distinctions are essential.

\emph{Descriptor versus proof.} Implicit patch and charted outputs are
compact descriptors of polynomial length. A descriptor denotes a unique
point because the restricted objective is strongly convex on the patch;
that point is a global optimizer because of a separate \emph{global proof
record}, which contains the pruning trace and the certificate checks. The
descriptor alone does not prove global optimality. The proof record can be
large on an individual draw; its expected size obeys the same bound as the
expected work.

\emph{Evaluation.} Given a descriptor and an integer $q\ge0$, a
\emph{$q$-bit evaluation} returns a rational point within Euclidean distance $2^{-q}$
of the optimizer and a rational interval of width at most $2^{-q}$ that
contains the optimal value. For patch outputs this costs $\poly_d(I+q)$ bit
operations: convex optimization in the bit model gives value accuracy, and
the verified modulus $g_0$ converts value accuracy into distance. On box,
polytope, and explicit-graph domains the evaluation also returns a feasible
rational point with a certified objective gap. An implicit graph need not
contain any rational feasible point: the equation $y^2=2$ with $y\in[1,2]$
is an example. For charted outputs on implicit graphs, an exactly feasible
point consists of rational free coordinates together with the algebraic
dependent coordinates that the chart assigns to them; rational
approximations of all coordinates need not be feasible. Evaluation does not
decide exact comparisons of the optimal value with a rational threshold,
and it does not identify the exact active set.
The exact-arithmetic companion shows why these are different requests:
exact active-coordinate recognition for a uniformly strongly convex cubic
can decide Square Root Sum while point approximation remains polynomial
\citep{companion-exact-arithmetic}.
If the patch has no free coordinate, its unique point is evaluated directly;
on an implicit graph this includes evaluation of its exact dependent lift.
No full-dimensional convex optimization routine is invoked on a point.

\emph{Algebraic sizes.} An algebraic output fixes one optimizer: all its
coordinates are functions of one root $\theta$, so coordinates of different
optimizers cannot be mixed. Its defining polynomial need not be minimal.
Algebraic outputs from fallback draws have length and $q$-bit refinement cost
at most $B\poly_d(I+b+q)$, where $B=2^{\poly_d(I)}$ is base-only; their
expected contribution is polynomial in the sampled bit length and requested
precision under the stated law, with any stated sampling-parameter factor
retained. A component output does not include a
defining polynomial of the total value, whose degree can grow exponentially
with the number of components, and it does not provide exact sign or
equality tests for sums of unrelated algebraic numbers. Rational
enclosures of the total value are obtained by refining each component.

\emph{Selection.} Unless a theorem says so, the returned optimizer is the
one the algorithm happens to certify. No canonical selection rule, such as
the lexicographically least or minimum-norm optimizer, is promised.

\subsection{Fixed-parameter and fixed-width bounds}\label{sec:model:parameters}

A bound is \emph{fixed-parameter} in a parameter $\kappa$ if it has the form
$f(\kappa)\cdot\poly(I)$, where the exponent of the polynomial does not
depend on $\kappa$. The low-negative-inertia, core, and integer-dimension
results are of this form, with numerical ratios such as
$\nu\diam(X)/\sigma$ or $L/\sigma$ included in $\kappa$. A bound is
\emph{fixed-width polynomial} if it is polynomial for each fixed value of a
parameter $p$ but its degree grows with $p$, for instance
$[nLw_{\max}/\sigma]^{O(p)}\poly(I)$. The sparse results are of this kind,
and they are not fixed-parameter bounds in the width. The limitation is
genuine for the method used: \cref{thm:lim:width} exhibits instances on
which its retained state count is $(\Omega(n/p))^{p/2}$ on every draw.

A bound that is polynomial in a numerical ratio is polynomial in its
magnitude, which can be exponential in its encoding length. The phrase
\emph{expected polynomial time} is used only when the relevant ratios are
polynomially bounded in $I$ and the structural parameters are fixed.
FPT means joint parameterization by structure and the displayed numerical
ratios, or parameterization by structure for each fixed bounded ratio.
If a ratio $R$ grows polynomially in $I$, a bound
$(1+R)^{O(k)}\poly(I)$ is $I^{O(k)}$, rather than FPT in $k$ alone.

\subsection{The original objective}\label{sec:model:original}

The theorems optimize the sampled objective $F_\gamma$. An exact optimizer
of $F_\gamma$ need not optimize $F_0$, and nothing is claimed about its
distance from the optimal set of $F_0$. The following elementary bound states
what a sampled solution certifies for the original problem. For compact $X$
and $\gamma\in\R^n$, let
\[
 \omega_X(\gamma)=\max_{x\in X}\gamma^{\mathsf T}x-\min_{x\in X}\gamma^{\mathsf T}x
\]
be the width of $X$ in the direction $\gamma$.

\begin{proposition}[Certified original-objective regret]\label{prop:model:regret}
Let $X$ be compact, let $F_0$ be continuous, and let $y\in X$ and $a\in\R$
satisfy
\[
 a\le\min_XF_\gamma\le F_\gamma(y)\le a+\delta.
\]
Then the interval with endpoints
\[
 a-\max_{x\in X}\gamma^{\mathsf T}x\qquad\text{and}\qquad F_0(y)
\]
contains $\min_XF_0$, and its length is at most $\delta+\omega_X(\gamma)$. In
particular $F_0(y)-\min_XF_0\le\delta+\omega_X(\gamma)$.
\end{proposition}

\begin{proof}
For every $x\in X$, $F_0(x)=F_\gamma(x)-\gamma^{\mathsf T}x\ge
a-\max_X\gamma^{\mathsf T}x$, which gives the lower endpoint. The upper
endpoint is the value of the feasible point $y$. Finally,
\[
 F_0(y)=F_\gamma(y)-\gamma^{\mathsf T}y\le a+\delta-\min_X\gamma^{\mathsf T}x
 =\Bigl(a-\max_X\gamma^{\mathsf T}x\Bigr)+\delta+\omega_X(\gamma).\qedhere
\]
\end{proof}

For an exact sampled optimizer $\hat x$ one may take $y=\hat x$, $\delta=0$,
and $a=\min_XF_\gamma$. The endpoints are then real numbers in the format of
the output; for implicit, charted, algebraic, and component outputs they
are generally irrational. Rational certified endpoints come from a
$q$-bit evaluation: a feasible point $y$ with a rational upper value, and a
rational lower bound $a$, whose distance from each other is the evaluation
error, which is added to $\delta$. The point $y$ must be exactly feasible.
On implicit graphs it consists of rational free coordinates and their exact
algebraic lift, and $F_0(y)$ is replaced by the upper end of a certified
rational enclosure; a coordinatewise rational approximation of the lift may
be infeasible and cannot be used.

If every perturbed coefficient satisfies $\abs{\gamma_i}\le\bar\sigma$, then
$\omega_X(\gamma)\le\bar\sigma\sum_iw_i$ on a mixed box, where the sum is
over the perturbed coordinates; under core-only noise on a unit core box this
is $k\bar\sigma$. For aligned noise $\omega_X(\gamma)\le\bar\sigma
\sum_{i=1}^k\omega_X(T_{i\cdot})$, where $T_{i\cdot}$ is the $i$th row of
$T$. Here $\bar\sigma$ bounds the largest magnitude in the support of the
marginal law: $\bar\sigma=\sigma$ for grid laws and
$\bar\sigma=(b+20)\sigma$ for Gaussian-like laws. On a box,
$\max_X\gamma^{\mathsf T}x$ is an endpoint calculation; on a polytope it is a
linear program. On a mixed polytope one may maximize over its continuous
relaxation, which keeps the lower endpoint valid and replaces $\omega_X$ by
the width of the relaxation.

For calibration, use widths in the coordinates of the chosen perturbation
model. Write $W_{\rm noise}=\sum_{i\text{ perturbed}}w_i$ for ambient or
core-only noise, and $W_{\rm noise}=\sum_i\omega_X(T_{i\cdot})$ for aligned
noise; polyhedral relaxations may provide upper bounds for these widths.
With row-specific aligned scales the bound is instead
$\sum_i\sigma_i\omega_X(T_{i\cdot})$. If $W_{\rm noise}=0$, the perturbation
is constant on $X$ and only the evaluation error needs to be charged.

For a route valid at every positive noise scale and $W_{\rm noise}>0$, let
$\varepsilon>0$, ensure $\bar\sigma W_{\rm noise}\le\varepsilon/2$, and
evaluate with total certificate error at most $\varepsilon/2$. This gives an
exactly feasible point and a rational certificate of original-objective gap
at most $\varepsilon$ on every draw. For grid laws take
$\sigma=\varepsilon/(2W_{\rm noise})$.
For Gaussian-like laws the support bound is exactly
$\bar\sigma=(b+20)\sigma$, with $b$ chosen by the effective base-data
schedule. Put $R=\max\{1,W_{\rm noise}/\varepsilon\}$ and let $I_0$ be the
length of the supplied data and accuracy, excluding the noise scale. There
is an effective nondecreasing polynomial $P$ such that the required accuracy
at $\sigma=2^{-t}$ satisfies $b(t)\le P(I_0+t)$, after absorbing fixed
encoding constants. Choose the least integer $t\ge0$ satisfying
\[
 2R\bigl(P(I_0+t)+20\bigr)\le2^t .
\]
Exponential growth dominates this fixed polynomial, so
$t=O(I_0+\log_2R+1)$. If $t\ge1$, failure at $t-1$ gives
$2^t<4R(P(I_0+t)+20)$; if $t=0$, $2^t=1\le R$. Consequently
\[
 \frac1\sigma\le R\,\poly(I_0+\log_2R+1),\qquad
 (b(t)+20)\sigma W_{\rm noise}\le\varepsilon/2 .
\]
This preserves polynomial sampling precision and includes the full support
factor in every numerical work bound.

Substitute this scale into the theorem's displayed bound, retaining every
other numerical factor. For \cref{thm:sp:main} under grid noise, for example,
the bound is
$C_0^p[4+(1+n/2)Lw_{\max}W_{\rm noise}/\varepsilon]^p\poly_d(I)$.
The work is polynomial in $1/\varepsilon$ for fixed structure and fixed
other numerical parameters; it is polynomial in $I$ and $1/\varepsilon$
only when those numerical factors are polynomially bounded. Binary-encoded
curvature, widths, and coupling-derived bounds may be exponentially large.
The strong-noise results retain the width bound but do not give this
prescribed-accuracy consequence: their lower noise-scale condition need not
permit the calibration. None of these guarantees yields exact optimization
of $F_0$ by taking a positive noise scale to zero.
This is consistent with the worst-case hardness of the unperturbed problems
discussed in \cref{sec:intro,sec:lim}.

\subsection{Why a finite law forces an exact treatment of ties}\label{sec:model:atom}

Under a continuous noise law, ties among optimizers and other degenerate
configurations typically have probability zero. Under a finite law they can
have positive probability that no refinement of the search removes. The
following example is the smallest instance of this phenomenon.

\begin{example}[A tie with probability $1/M$]\label{ex:model:tie}
Let $X=[0,1]^2$, $F_0(x)=2x_1x_2-x_1-x_2$, and let $\gamma_1,\gamma_2$ be
independent with law $U_{\sigma,M}$, where $0<\sigma\le1/4$. The event
$\gamma_1=\gamma_2$ has probability exactly $M\cdot M^{-2}=1/M$.

On this event, put $s=1-\gamma_1\in[3/4,5/4]$. The minimum value of
$F_\gamma$ is $-s$, and
\[
 F_\gamma(x)+s=s(1-x_1-x_2)+2x_1x_2 .
\]
If $x_1+x_2\le1$, both terms are nonnegative, and they vanish together only
at $(1,0)$ and $(0,1)$. If $x_1+x_2>1$, then $x_1x_2>0$, and
$(1-x_1)(1-x_2)\ge0$ gives $x_1+x_2-1\le x_1x_2$, so
$F_\gamma(x)+s\ge(2-s)x_1x_2>0$. Hence $F_\gamma$ has exactly the two
global minimizers $(1,0)$ and $(0,1)$ on every draw of this event.

Consequently no certificate that exhibits a strongly convex patch containing
every global optimizer can succeed on this event, at any refinement depth:
a strictly convex function has at most one minimizer on a convex set. The
point-growth modulus is zero, so any tail bound of the form
$\Pr\{g_*<\varepsilon\}\le c\varepsilon+\beta$ for the grid law must have
$\beta\ge1/M$ for all $\varepsilon>0$. An algorithm that is correct on every
draw must solve this event by other means. The closure-based results do so
with their same-draw fallback, whose expected contribution on this event is
at most $M^{-1}$ times its cost; choosing $M$ large compared with the
base-only factor $B$ keeps this contribution polynomial. Under any
absolutely continuous law for $(\gamma_1,\gamma_2)$ the tie has probability
zero, which is why continuous-noise analyses can ignore it.
\end{example}

The example shows why the probability estimates of the later sections have
the form ``continuous bound plus atomic term'', and why the grid size $M$ is
chosen only after the fallback factor and the stopping depth are known.
\Cref{ex:count:atom} shows the companion effect on counts: a fixed atom can
keep the number of retained cells growing with the refinement level.
